% =========================================================
% Proof Stub: Convergent Series Has Cauchy Partial Sums
% Source: volume-iii/analysis/sequences/notes/sequences/notes-cauchy.tex
% Mode: standard
% =========================================================

\clearpage
\phantomsection
\hypertarget{proof-RA-ABB-C-conv-series-cauchy}{}

\subsubsection[Convergent Series Has Cauchy Partial Sums]{Proof --- Convergent Series Has Cauchy Partial Sums}

\bigskip

\noindent
\textbf{Source.}
Abbott, \emph{Understanding Analysis}; see \texttt{notes-cauchy.tex}.

\vspace{0.75em}

\noindent
\textbf{Goal.}
If $\sum a_n$ converges, then its partial sums $(s_n)$ form a Cauchy sequence.

\vspace{0.75em}

\noindent
\textbf{Logical form.}
\[
  \sum a_n \text{ converges} \Rightarrow \forall \varepsilon > 0,\; \exists N,\; m > n \ge N \Rightarrow \left|\sum_{k=n+1}^m a_k\right| < \varepsilon.
\]

\vspace{0.75em}

\noindent
\textbf{Proof strategy.}
$(s_n)$ converges (by hypothesis), so it is Cauchy by the theorem ``every convergent sequence is Cauchy.'' The partial-sum form just unpacks $|s_m - s_n|$.

\vspace{0.75em}

\noindent
\textbf{Proof.}
\begin{proof}
\Claim If $\sum a_n$ converges, then its partial sums $(s_n)$ form a Cauchy sequence.

\Given 

\Goal 

\Strategy $(s_n)$ converges (by hypothesis), so it is Cauchy by the theorem ``every convergent sequence is Cauchy.'' The partial-sum form just unpacks $|s_m - s_n|$.

\medskip

% --- s_n → S (hypothesis) ---
% --- by Convergent ⇒ Cauchy: (s_n) is Cauchy ---
% --- unpack: |s_m - s_n| = |∑_{k=n+1}^m a_k| ---

\end{proof}

\begin{remark*}[Proof shape]
Direct corollary: convergent $\Rightarrow$ Cauchy, applied to partial sums.
\end{remark*}

\begin{remark*}[Dependencies]
\begin{itemize}
  \item Every convergent sequence is Cauchy.
  \item Definition of partial sums: $s_m - s_n = \sum_{k=n+1}^m a_k$.
\end{itemize}
\end{remark*}
